\section{什么是 Agentic Memory？}

Agentic Memory 不是单一组件，而是一个在后台运作的\textbf{完整系统}——包括不同类型的存储、信息检索方式以及智能管理策略，使 Agent 能够跨时间维持上下文。

记忆系统同时承担三个截然不同的职责：

\begin{enumerate}[leftmargin=2em]
    \item \textbf{Continuity（连续性）}——关乎身份认同。Agent 知道你是谁、你的偏好是什么、你们之前一起构建了什么。没有连续性，每次交互都像从头开始。
    \item \textbf{Context（上下文）}——关乎当前任务。刚刚发生了什么、调用了哪个工具、返回了什么结果、下一步该做什么。它是多步工作流不崩溃的保障。
    \item \textbf{Learning（学习）}——关乎持续进步。理解什么有效、什么无效，并在长期中逐步改进决策，而非重复犯同样的错误。
\end{enumerate}

\begin{knowledgebox}{三大职责的协同}
一个设计良好的 Agent 记忆系统同时处理以上三个方面，为每个方面使用不同的存储后端。三者协同让 Agent 表现得一致、可靠，且在每次交互后都变得更聪明。
\end{knowledgebox}

\subsection{本章小结}

Agentic Memory 是一个多层系统，同时服务于身份连续性、任务上下文和经验学习三个目标。

